% =====================================================================
%  CHAPTER 10: 儿童少年心理卫生问题及精神障碍
%  参考PPT：第十章 儿童少年心理卫生问题及精神障碍（龚雯洁，中南大学）
% =====================================================================
\chapter{儿童少年心理卫生问题及精神障碍}
\setcounter{chapter}{10}
\chapterintro{
\textbf{掌握：}儿童少年心理健康的基本特征、影响因素；一般心理卫生问题（情绪问题、行为问题、适应问题）的表现及干预；常见精神障碍（ADHD、ASD、TD、SLD、焦虑障碍、抑郁障碍、双相障碍、OCD、品行障碍、神经性厌食症）的主要表现及干预方法。\newline
\textbf{熟悉：}心理卫生与精神障碍的区别；精神障碍的流行特征。
}

% ----- 10.1 儿童少年心理健康及其影响因素 -----
\section{儿童少年心理健康及其影响因素}

\subsection{儿童心理健康的基本特征}

\begin{defbox}
\textbf{儿童少年心理健康（child and adolescent mental health）：}指儿童少年在情感、行为、社会和认知领域达到特定年龄阶段的要求。对其的认知和理解包括以下几个方面：

\begin{enumerate}[leftmargin=*]
    \item \textbf{相对性}——心理健康是一种相对的状态，而非绝对的"完美"。
    \item \textbf{持续性}——短暂的情绪波动或轻微行为异常，若能迅速自愈，不应视为不健康。
    \item \textbf{可描述性}——心理健康可以通过一系列具体标准来描述（如情绪稳定、适应环境、有效应对压力和良好社会行为等），但并非每个个体都能完全满足标准。
    \item \textbf{多元性}——心理健康不仅仅是单一维度，而是生物、心理与社会因素共同作用的结果。
    \item \textbf{发展性}——心理健康超越了单纯的"无病状态"，还涵盖了儿童少年幸福感的体验和成长发展的潜力。
\end{enumerate}
\end{defbox}

\begin{keybox}
\textbf{儿童少年心理健康的基本特征：}

尽管健康的心理没有"标准模板"，但心理健康的儿童少年有其基本特征：

\begin{enumerate}[leftmargin=*]
    \item \imp{智力发展正常}——智力是衡量一个人心理健康的重要标志之一，智力发展水平要符合生理年龄。

    \item \imp{情绪稳定且反应适度}——心理活动和行为模式应和谐统一，对外部刺激反应适度，表现既不异常敏感也不异常迟钝，情绪波动在合理的范围内，并具有一定应变和应对能力。

    \item \imp{心理行为特点与年龄相符}——具有与生理年龄相符的心理和行为特征及模式。

    \item \imp{人际关系的心理适应}——心理健康的儿童少年，有积极、良好的人际关系，在集体中能与人和睦相处，悦纳自己，认同他人。

    \item \imp{个性稳定和健全}——表现出健康的精神风貌，客观积极的自我意识，行为符合社会道德规范，能适度耐受各种压力和应激。
\end{enumerate}
\end{keybox}

\subsection{儿童少年心理健康的影响因素}

\begin{keybox}
\textbf{影响儿童少年心理健康的三类因素：}

儿童少年心理健康受生物因素、心理因素和社会因素的综合影响。每一类因素均包含危险因素（增加心理健康问题风险）和保护因素（降低风险、促进心理健康）。
\end{keybox}

\subsubsection{生物因素}

\begin{defbox}
\textbf{一、危险因素：}

\begin{enumerate}[leftmargin=*]
    \item \imp{遗传易感性}——个体因遗传基因而对某些心理健康障碍具有较高风险。
    \item \imp{神经生物学异常}——大脑不同区域与情绪调节、认知功能和行为控制密切相关。神经递质的失衡可能导致情绪和行为问题。
    \item \imp{躯体健康状况不良}——伴随的身体健康问题会导致儿童少年感到痛苦和焦虑，影响情绪和行为。
    \item \imp{智力低下}——智力低下的儿童少年容易面临情绪、行为问题以及社交适应困难。这些问题通常源于认知功能受限、社交技能不足及社会环境的挑战，可能导致自尊心低下、情绪困扰和生活适应障碍。
    \item \imp{母孕期不良因素}——不良的孕期健康状况可能导致胎儿发育不良，增加出生后认知和情绪问题的风险。孕期心理健康问题也可能对胎儿的神经生物学基础造成潜在影响，进而影响儿童少年的情感和社交发展。
\end{enumerate}

\textbf{二、保护因素：}

\begin{enumerate}[leftmargin=*]
    \item \imp{良好的躯体健康}——躯体健康直接影响情绪、认知和行为。均衡的营养、规律的锻炼和良好的睡眠习惯促进躯体功能正常，增强免疫力，减少因躯体疾病引发的心理健康问题。良好的躯体健康为儿童少年提供稳定的生理基础，使其能够更好地应对生活中的挑战。
    \item \imp{适龄的发育水平}——儿童少年期的躯体与智力发育与心理发展密切相关。正常的发育水平有助于儿童形成健康的自我形象，增强自尊心和自信心，从而积极参与各类活动，促进良好的人际关系和心理适应能力。
\end{enumerate}
\end{defbox}

\subsubsection{心理因素}

\begin{defbox}
\textbf{一、危险因素：}

\begin{enumerate}[leftmargin=*]
    \item \imp{消极认知模式}——个体在儿童少年时期形成的思维方式对情绪和社交能力有重要影响。
    \item \imp{情绪调节能力不足}——缺乏有效的情绪调节策略的儿童少年，可能在面对压力或挫折时表现出冲动行为或情绪崩溃，长期的调节困难可能加重心理健康问题。
    \item \imp{"难养型"气质}——具有"难养型"气质的儿童少年在生理节律和情绪反应上表现不规律，适应新环境和变化的能力较差，容易受到负面情绪影响，增加焦虑和抑郁的风险。
    \item \imp{自卑/自负}——自卑的儿童少年常自我评价低，过于关注他人评价，情绪低落且回避竞争；而自负的儿童少年则表现出自我中心、过度自夸，缺乏客观的自我与他人关系评估。
\end{enumerate}

\textbf{二、保护因素：}

\begin{enumerate}[leftmargin=*]
    \item \imp{良好的自我意识}——自我意识是个体对自我的认识和态度，促进社会化和完善人格特征。良好的自我意识使儿童少年能够清晰地理解自身的情感、需求和价值观。
    \item \imp{利于适应的人格特征}——适应性强的人格特征通常包括开放性、责任感和情绪稳定性。
    \item \imp{牢固的人际关系}——家庭、朋友和学校构成的支持网络提供情感支持和归属感，降低孤独感和焦虑感。
\end{enumerate}
\end{defbox}

\subsubsection{社会因素}

\begin{defbox}
\textbf{一、危险因素：}

\begin{enumerate}[leftmargin=*]
    \item \imp{家庭危机}——家庭是儿童成长的首要环境。家庭功能不全会导致缺乏安全感和情感支持，增加焦虑和抑郁的风险。
    \item \imp{学业压力与校园暴力}——学业压力，尤其是在应试教育下，可能导致焦虑和抑郁情绪加重。
    \item \imp{社区环境的不稳定}——城市化、移民潮和社区环境的混乱可能使儿童暴露于恐惧和不安之中。歧视和边缘化增加心理压力。网络的普及使儿童易接触到不良内容，导致不良行为和心理健康问题的发生。缺乏有效的社区支持限制了儿童少年与同伴和成人建立健康关系的可能性，进一步加剧孤独感和无助感。
\end{enumerate}

\textbf{二、保护因素：}

\begin{enumerate}[leftmargin=*]
    \item \imp{家庭关怀}——温暖和支持性的家庭关系能有效降低儿童少年心理健康问题的风险。
    \item \imp{支持性的学校环境}——关怀、支持和激励的学校氛围为学生提供情感支持，促进心理韧性。良好的学业成就带来正性强化，有助于树立自我价值感。
    \item \imp{社区联系与文化包容性}——与社区及其组织的紧密联系增强儿童少年的归属感和社会支持，从而降低心理健康风险。
\end{enumerate}
\end{defbox}

\subsubsection{危险因素与保护因素汇总}

\begin{keybox}
\textbf{儿童少年心理卫生的危险因素和保护因素汇总表：}

\begin{table}[H]
\centering
\caption{儿童少年心理卫生的危险因素和保护因素}
\small
\begin{tabularx}{\textwidth}{l>{\raggedright\arraybackslash}p{5cm}>{\raggedright\arraybackslash}p{5cm}}
\hline
\textbf{层次} & \textbf{危险因素} & \textbf{保护因素} \\
\hline
\imp{生物因素} &
遗传易感性；智力低下；\newline
神经生物学异常；\newline
母孕期不良因素；\newline
躯体健康状况不良 &
良好的躯体健康；\newline
适龄的躯体发育 \\
\hline
\imp{心理因素} &
消极认知模式；\newline
"难养型"气质；\newline
情绪调节能力不足；\newline
自卑/自负 &
从经验中学习的能力；\newline
解决问题能力强；\newline
良好的自我意识；\newline
适应良好的人格特征 \\
\hline
\imp{社会因素（家庭）} &
更换照料者；家庭管理凌乱、冲突；\newline
缺乏家庭支持；家庭关系破裂；\newline
不良的家庭教育；贫穷 &
良好的物质生活条件；家庭成员健全；\newline
支持性的家庭关系；\newline
有机会积极参与家庭生活；\newline
明确的家庭规范 \\
\hline
\imp{社会因素（学校）} &
学业表现差；不适当的教育；\newline
学校未提供适当的环境；校园欺凌 &
关怀、支持、保护和激励的学校环境；\newline
与同伴间的积极关系；\newline
有参与学校生活的机会；\newline
从学业成就中得到正性强化 \\
\hline
\imp{社会因素（社区）} &
社会转型；暴露于暴力；\newline
混乱的社区环境；\newline
暴露于网络不良内容；\newline
歧视和边缘化；缺乏"归属"感 &
与社区、社区组织保持联系；\newline
文化包容和多样性；\newline
有机会参与有益的业余活动；\newline
安全的社会环境；\newline
有组织的支持性社区 \\
\hline
\end{tabularx}
\end{table}
\end{keybox}

% ----- 10.2 儿童少年一般心理卫生问题 -----
\section{儿童少年一般心理卫生问题}

\subsection{概述}

\begin{defbox}
\textbf{心理卫生问题（mental health problems）：}心理健康方面存在的偏倚统称为心理卫生问题。2019年的调查报告显示，目前全球20\%左右的10至19岁儿童少年存在心理卫生问题，且这些问题导致了约16\%的疾病和伤害。

2021—2025年联合国儿童基金会（UNICEF）估算，至少有3000万17岁以下的儿童少年面临情绪或行为问题。

儿童少年的一般心理卫生问题可从情绪、行为和适应等方面了解。与成人的心理健康问题不同的是，儿童少年更易受环境的影响，他们的困境往往反映了其家庭与学校等紧密环境的养育、教育问题。重大的社会环境变化也会加剧儿童少年群体的心理问题，因此在疫情后对这一领域的关注尤为重要。
\end{defbox}

\subsection{情绪问题}

\begin{keybox}
\textbf{情绪问题的定义与分类：}

儿童少年情绪问题主要包括\imp{焦虑}、\imp{抑郁}和\imp{孤独感}，严重时可发展为焦虑障碍或心境障碍，甚至引发自伤或自杀行为。

\begin{itemize}[leftmargin=*]
    \item \textbf{焦虑}——通常源于对令人恐惧或不确定情况的担忧。
    \item \textbf{抑郁}——是对生活中压力事件（如家庭变故或适应困难）的反应。
    \item \textbf{孤独感}——指个体在主观上感受到与社会隔绝、孤立无援的心理状态。儿童时期的孤独感可能预测成年后的孤独体验，有孤独感的儿童可能会错过许多与同龄人交流和学习重要终身技能的机会。
\end{itemize}
\end{keybox}

\begin{defbox}
\textbf{情绪问题的具体表现：}

情绪问题的界定较为困难，可以表现为情绪不稳定、紧张焦虑、抑郁、强迫行为、过度任性、冲动、暴躁易怒、沮丧、胆小退缩和恐惧发作等，也可有啃咬指甲等行为表现。

\begin{enumerate}[leftmargin=*]
    \item \imp{焦虑的表现}——常表现为紧张、不安或过度担忧的情绪，尤其在面临考试或社交场合时。他们往往会回避特定情境，例如拒绝参加社交活动或逃避学校，可能出现心跳加速、呼吸急促和出汗、入睡困难或频繁做噩梦等生理症状。

    \item \imp{抑郁的表现}——常表现为悲伤、沮丧和失望，伴随无助感和绝望感。在行为上，抑郁的儿童少年常出现兴趣减退、社交退缩和学习动机降低等，生理反应包括疲惫感、食欲变化和睡眠质量下降。

    \item \imp{孤独的表现}——常表现为较强的依恋需求，对父母的陪伴有较高的渴望。
\end{enumerate}
\end{defbox}

\begin{tipbox}
\textbf{缓解不良情绪的干预措施：}

\begin{enumerate}[leftmargin=*]
    \item \imp{识别并解除环境中的压力源}，鼓励儿童少年开展放松练习和活动，如增加身体活动、阅读、深呼吸练习、改善饮食习惯等。
    \item \imp{促进同伴关系}以及加强学校和家庭的支持，帮助他们控制和减轻情绪负担，增强心理韧性。
    \item 若儿童少年的情绪持续时间长且程度加剧时，可\imp{寻求专业心理援助人员}的帮助。
\end{enumerate}
\end{tipbox}

\subsection{行为问题}

\begin{defbox}
\textbf{行为问题（behavioral problems）：}通常指的是在日常生活中出现的不符合其年龄或社会期望的行为，其成因复杂且多样，尚未明确单一的诱因。

儿童少年倾向于通过行为表达情绪，这往往是他们在适应环境或表达内心冲突时的反应。由于儿童少年可能缺乏成熟的应对机制，相对于成年人通常可以运用更复杂的应对策略，他们更容易表现出冲动行为。
\end{defbox}

\begin{keybox}
\textbf{儿童少年常见行为问题的分类：}

\subsubsection*{一、暴力行为}

暴力行为是指蓄意运用躯体力量或权力对自身、他人、群体或社会进行威胁或伤害的行为。

\subsubsection*{二、非自杀性自伤（NSSI）}

非自杀性自伤可以认为是针对自己的暴力行为，指个体在没有自杀意图的情况下，故意直接破坏自身身体组织的行为，而且这种行为通常不被当前社会文化所认可。

\begin{itemize}[leftmargin=*]
    \item 全球范围内儿童少年NSSI的终身发生率为22.1\%，12个月发生率为19.5\%。
\end{itemize}

\subsubsection*{三、物质使用}

主要包括酒精、烟草、非法药物和吸入剂等。

\begin{itemize}[leftmargin=*]
    \item 全球18岁以下儿童少年中，酒精使用率为41.7\%，烟草15.6\%，非法药物10.1\%，大麻6.3\%。
    \item 在中国，6$\sim$22岁学生中，曾尝试吸烟和饮酒的比例分别为9.3\%和26.6\%。
\end{itemize}
\end{keybox}

\begin{tipbox}
\textbf{行为问题的干预措施：}

\begin{enumerate}[leftmargin=*]
    \item \imp{个体层面}——需关注行为问题出现的具体原因，从病因出发针对儿童少年开展个性化的心理干预措施，包括认知行为疗法、动机访谈等心理疗法。
    \item \imp{环境层面}——针对儿童青少年行为问题的干预措施也应综合考量家庭、学校等多方面的环境因素，并重视社会支持系统的建立健全。
\end{enumerate}

\vspace{4pt}
\textbf{INSPIRE七项策略：}为预防儿童少年暴力，世界卫生组织（WHO）提出"消除针对儿童的暴力行为的七项策略"（INSPIRE），在个人、亲密关系、社区、社会四个层面上系统处理风险因素和保护因素。
\end{tipbox}

\subsection{适应问题}

\begin{defbox}
\textbf{适应问题（adjustment problem）：}指在经历生活变化或压力事件时，出现的暂时性情绪、行为或学习上的困扰，如悲伤、绝望、非自杀性自伤、学习困难等，常伴有头痛、腹痛等身体症状。这些问题反映了个体在面对环境变化或心理压力时所遇到的挑战，通常在成长过程中可以逐渐克服。
\end{defbox}

\begin{keybox}
\textbf{适应问题的分类与表现：}

适应问题可以分为社交、学业和家庭适应问题，具体表现包括：

\begin{itemize}[leftmargin=*]
    \item 与同龄人互动困难、无法建立和维持友谊、没有朋友、孤独、被排斥或欺凌。
    \item 缺乏同理心和社交技巧、在社交场合中感到极度不安或恐惧、可能会避免参加社交活动。
    \item \imp{升学相关压力}——由于升学对儿童少年来说是重要的压力事件，在进入小学（6岁）和小学升入中学（12岁）阶段是他们因头痛、腹痛等身体不适寻求公共卫生服务的高峰期。
\end{itemize}
\end{keybox}

\begin{tipbox}
\textbf{适应问题的干预措施：}

虽然引发适应问题的原因复杂多样，但儿童少年仍可以采取以下措施减轻压力，更好地适应变化：

\begin{enumerate}[leftmargin=*]
    \item \imp{建立支持系统}——形成由家人、朋友、同龄人或社会团体组成的网络，在面临挑战、压力或困难时提供支持和鼓励。
    \item \imp{练习自我管理}——当心烦意乱或感受到压力时，儿童少年可通过洗热水澡、读书、写日记、出去散步或其他喜欢的活动放松身心。
    \item \imp{保持健康的生活方式}——鼓励其健康饮食，规律运动，保持身心健康。
\end{enumerate}
\end{tipbox}

% ----- 10.3 儿童少年精神障碍 -----
\section{儿童少年精神障碍}

\subsection{概述}

\begin{defbox}
\textbf{儿童少年精神障碍（mental disorders）：}指个体在认知、情绪调节或行为上的严重紊乱，通常反映了心理、生物或发展过程中的功能失调。
\end{defbox}

\begin{keybox}
\textbf{儿童少年精神障碍的三大特征：}

\begin{enumerate}[leftmargin=*]
    \item \imp{个体的痛苦体验}——经历显著的情感痛苦，这种痛苦可能使他们在日常生活中感到无能为力或绝望。情感上的困扰不仅影响心理健康，还可能引发躯体症状，形成身体与心理之间的恶性循环。

    \item \imp{功能损害}——精神障碍会导致个体在多个领域的功能受损，包括情感调节、社交互动和学业表现。认知功能的下降，可能直接影响他们的学业和自信心，进一步加剧社交隔离感。

    \item \imp{未及时处理后恶化}——如果精神障碍未能及时得到干预和治疗，个体的症状可能会加重，导致功能损害进一步加深，甚至引发身心痛苦、伤残或极端情况下的自杀风险。
\end{enumerate}
\end{keybox}

\begin{exambox}
\textbf{儿童少年精神障碍的流行情况：}

\begin{itemize}[leftmargin=*]
    \item 精神障碍的发病高峰年龄为14.5岁，63$\sim$75\%的精神障碍在25岁之前发病。
    \item 由精神障碍造成的伤残调整生命年在儿童少年时期明显上升，在25$\sim$34岁达到峰值。
    \item 全球范围内，有8.8\%的儿童少年被诊断患有精神障碍。
    \item 中国2020年对6$\sim$16岁在校学生进行的流行病学调查指出，精神障碍的总患病率为17.5\%。
\end{itemize}
\end{exambox}

\subsection{神经发育障碍}

\begin{defbox}
\textbf{神经发育障碍：}指儿童少年时期，由于大脑发育异常或延迟，导致学习、运动、社交等技能的显著困难。
\end{defbox}

\subsubsection{注意缺陷多动障碍（ADHD）}

\begin{keybox}
\textbf{注意缺陷多动障碍（attention-deficit/hyperactivity disorder, ADHD）：}也称多动症，是以注意力不集中、多动、冲动为特征的综合征。

\vspace{4pt}
\textbf{（1）流行特征：}

\begin{itemize}[leftmargin=*]
    \item 全球范围内3$\sim$12岁儿童少年中ADHD的现患率为7.6\%，13$\sim$18岁少年中现患率为5.6\%。
    \item 2020年流行病学调查显示，中国儿童少年ADHD患病率为6.4\%，约2300万人。
    \item 60\%$\sim$80\%患儿的ADHD可持续到青少年期，50.9\%持续为成人ADHD。
\end{itemize}

\vspace{4pt}
\textbf{（2）主要表现（按年龄阶段）：}

\begin{itemize}[leftmargin=*]
    \item \imp{学龄前期}——患儿容易转移注意力，似听非听；过分喧闹和捣乱，无法接受幼儿园教育；可能出现明显攻击行为，不好管理。
    \item \imp{学龄期}——患儿不能完成指定任务，容易转移注意力，不能集中精神；同时烦躁、坐立不安，有过多的语言；自制力差，难以等待按顺序做事情，言语轻率。
    \item \imp{青少年时期}——主要表现为不能完成作业，容易转移注意力；主观上有不安宁的感觉；自制力差，经常参与危险性活动。
\end{itemize}

\vspace{4pt}
\textbf{（3）预防及干预：}

ADHD的病因和发病机制尚不完全清楚，目前认为ADHD的发生是在胚胎期和婴儿早期由复杂的遗传易感性与暴露环境多种不利因素协同作用的结果，因此要对具有高危因素（如有患ADHD的兄弟姐妹、母亲孕期吸烟饮酒等）的儿童进行监测和早期识别。

ADHD的干预因年龄而异：
\begin{itemize}[leftmargin=*]
    \item 对于\imp{学龄前儿童}，首选非药物治疗，包括心理教育、心理行为治疗、特殊教育和功能训练，并开展家长培训和学校干预。
    \item \imp{6岁以后}采用药物治疗和非药物治疗相结合的综合治疗，以帮助患儿以较低用药剂量达到最佳疗效。
\end{itemize}
\end{keybox}

\subsubsection{孤独症谱系障碍（ASD）}

\begin{keybox}
\textbf{孤独症谱系障碍（autistic-spectrum disorder, ASD）：}是一种以社交沟通障碍和重复刻板行为为主要特征的发育障碍性疾病。

\vspace{4pt}
\textbf{（1）流行特征：}

\begin{itemize}[leftmargin=*]
    \item WHO的数据显示，全世界大约每100名儿童中就有1名患有ASD。
    \item 美国CDC的自闭症和发育障碍监测（ADDM）网络系统估计，2020年美国每36名8岁儿童中就有1名患有ASD，患病率约为2.76\%。
    \item ASD在男孩中的患病率约为女孩的4倍。
    \item 我国6$\sim$12岁学龄期儿童ASD的患病率约为0.7\%。
\end{itemize}

\vspace{4pt}
\textbf{（2）主要表现（核心症状）：}

\begin{enumerate}[leftmargin=*]
    \item \imp{社交交流障碍}——患儿难以参与到其他人中，难以分享兴趣、情感和进行非言语交流，难以发展、维持和理解人际关系。
    \item \imp{局限性、重复性的行为、兴趣或活动模式}——患儿表现为躯体刻板运动（拍手、弹手指）、重复使用物体（旋转硬币、摆放玩具）、重复言语（刻板使用单词、短语或韵律模式）、刻板地遵守日常惯例和行为模式等。
    \item \imp{感官反应过度或反应不足}——对光线和旋转物体特别感兴趣、对特定气味、味道或质地表示极度厌恶、对疼痛或温度感觉麻木。
\end{enumerate}

\vspace{4pt}
\textbf{（3）预防及干预：}

国家卫生健康委办公厅《0$\sim$6岁儿童孤独症筛查干预服务规范（试行）》要求乡镇卫生院、社区卫生服务中心等基层医疗卫生机构使用"儿童心理行为发育问题预警征象筛查表"对辖区0$\sim$6岁儿童实施11次初筛：
\begin{itemize}[leftmargin=*]
    \item 1岁以内婴儿期4次，分别在3、6、8、12月龄时。
    \item 1至3岁幼儿期4次，分别在18、24、30、36月龄时。
    \item 学龄前期3次，分别在4、5、6岁时。
\end{itemize}

干预的原则有\imp{早期}、\imp{个体化}、\imp{科学循证}、\imp{长程高强度}、\imp{基层为主}和\imp{家庭参与}，以改善社会交往、语言和非语言沟通能力为核心内容，以行为疗法为基本手段，结构化教育与自然情境下养育为干预基本框架，培养生活自理和独立生活能力，减少不适应行为，提高生存技能和交往能力。
\end{keybox}

\subsubsection{抽动障碍（TD）}

\begin{keybox}
\textbf{抽动障碍（tic disorders, TD）：}是一种起病于儿童期、以不自主的、突发的、快速的、反复的单一或多个部位的肌肉运动抽动和/或发声抽动为主要临床特征的神经发育障碍性疾病，通常共患各种精神和/或行为障碍，如ADHD、强迫性障碍、焦虑障碍、抑郁障碍和睡眠障碍等。

\vspace{4pt}
\textbf{（1）流行特征：}

\begin{itemize}[leftmargin=*]
    \item 国内报告6$\sim$16岁在校学生的TD患病率为2.3\%$\sim$3.0\%。
    \item TD对男孩的影响比女孩更大，比例为（1.5$\sim$4）:1。
\end{itemize}

\vspace{4pt}
\textbf{（2）主要表现：}

TD以\imp{运动性抽动}和\imp{发声性抽动}为临床核心症状。

\begin{itemize}[leftmargin=*]
    \item \textbf{运动抽动}——一般以眼肌、面肌和颈部肌肉抽动为多见，如反复眨眼、努嘴、缩鼻、摇头、耸肩、张口等。
    \item \textbf{发声抽动}——主要表现为刻板、重复地发出"哼"声，或反复地清嗓音、咳嗽等。
\end{itemize}

根据抽动的持续时间、参与的身体部分和肌肉群，运动抽动和发声抽动可再细分为简单性和复杂性：
\begin{itemize}[leftmargin=*]
    \item \textbf{简单性抽动}——包括单个肌肉或局部肌肉群的短暂收缩，表现为简单的运动或发声。
    \item \textbf{复杂性抽动}——会激活更多的肌肉群，表现为目标导向的或类似有目的的运动或单词或短语的发声。
\end{itemize}

40\%$\sim$55\%的患儿报告在运动抽动或发声之前有\imp{先兆性冲动}，表现为局部压、痒、痛、热或冷的感觉，或其他奇怪的感觉，也被称为感觉抽动。
\end{keybox}

\subsubsection{特定学习障碍（SLD）}

\begin{keybox}
\textbf{特定学习障碍（specific learning disorder, SLD）：}指在排除言语功能发育异常、获得性脑损伤、情绪以及环境等原因后，儿童少年在阅读书写、拼字、表达、计算等方面的基本心理过程存在一种或一种以上的特殊性障碍。

\vspace{4pt}
\textbf{（1）流行特征：}

\begin{itemize}[leftmargin=*]
    \item 累及阅读、书写和数学等学业领域的SLD的患病率在不同语言和文化的学龄儿童中为5\%$\sim$15\%。
    \item 阅读和书写障碍是最常被诊断出的SLD，数学障碍相对较少。
    \item 根据美国全国代表性数据，6$\sim$17岁儿童少年学习障碍患病率估计为8.83\%。
    \item 来自中国多地的调查结果显示，小学生阅读障碍患病率为3.9$\sim$9.7\%。
\end{itemize}

\vspace{4pt}
\textbf{（2）主要表现：}

\begin{enumerate}[leftmargin=*]
    \item \imp{阅读困难}——文字阅读正确率低，朗读时不准确或慢而吃力，大声读单个词时不正确或慢且犹豫，常常猜词，难以读出词。
    \item \imp{难以理解所读内容}——可正确地读出内容，但不能理解其顺序、关系、推论或更深层次的意思。
    \item \imp{拼写困难}——拼写准确性低，可能添加、省略或替代元音或辅音，常见"符号镜像"现象，如将p看成q，b看成d，m看成w，was看成saw，6读成9，"部"认作"陪"等。
    \item \imp{书面表达困难}——在句子里犯多种语法或标点符号错误，写字潦草难看，涂擦过多，段落组织凌乱，书面表达的意思不清。
    \item \imp{难以掌握数感、数字事实或计算}——对数、量及其关系理解差，常借助手指计数来计算个位数加法，而不是像同龄人那样回想数学事实，不能理解算术运算、可能转换步骤。
    \item \imp{数学推理困难}——运用数学概念、事实或步骤解决数量问题时存在严重困难。
\end{enumerate}

\vspace{4pt}
\textbf{（3）预防及干预：}

儿童早在3岁就可以表现出SLD迹象和指标，如发音问题、找词困难、难以押韵、难以遵循指示或学习常规等，教师可根据需要对儿童入学后进行筛选，每2$\sim$3年检查一次，以便于更早地识别SLD。

SLD通常采用综合矫治方法，根据儿童青少年的年龄、类型、严重程度、临床表现和心理评估结果制定个性化干预方案。常规程序包括\imp{个性化教育计划（individualized education program, IEP）}、个别指导计划和普通学校的特殊教育等。IEP以普通教学为基础，有明确的学年教育安置，定有相关的教育服务目标，是有效干预SLD的关键。
\end{keybox}

\subsection{焦虑障碍}

\begin{defbox}
\textbf{焦虑障碍（anxiety disorder）：}指无明显客观原因而出现的以不安和恐惧为主的情绪障碍，伴有明显的自主神经功能异常表现。2021年中国6$\sim$16岁在校学生焦虑障碍的患病率是4.7\%。
\end{defbox}

\subsubsection{广泛性焦虑障碍（GAD）}

\begin{keybox}
\textbf{广泛性焦虑障碍（generalized anxiety disorder, GAD）：}是最常见的焦虑障碍亚型，以持续、显著的紧张不安，伴有自主神经功能亢进和过分警觉为特征，并不局限于甚至并不特别突出地表现于任何一种特殊的环境或情况。

\vspace{4pt}
\textbf{（1）流行特征：}

\begin{itemize}[leftmargin=*]
    \item 中国首个全国性的儿童少年精神障碍流行病学调查显示，中国儿童少年GAD流行率为1.3\%。
    \item GAD在女性中比在男性中更频繁地发生，患儿的家庭成员中可能同时存在焦虑障碍。
\end{itemize}

\vspace{4pt}
\textbf{（2）主要表现：}

\begin{itemize}[leftmargin=*]
    \item 患儿常不切实际地担心过去和将来可能发生的各种事件，经常需要得到抚慰和保证，并且自我意识过强、多疑、完美主义。
    \item 儿童内心的焦虑不安不易被其养育者察觉，往往通过头痛、胃痛、失眠、心悸、口干等躯体症状表达出这种情绪。
    \item GAD在儿童期常合并分离焦虑障碍和社交焦虑障碍，少年期易合并抑郁障碍。
\end{itemize}

\vspace{4pt}
\textbf{（3）预防及干预：}

GAD的预防重点在于增强心理韧性和情绪调节能力，培养儿童有效的情绪表达与应对策略至关重要，父母和教师可以通过指导儿童识别和管理焦虑情绪，如使用放松训练、认知重构等方法，帮助其更好地应对压力情境。

对患儿的干预应遵循\imp{全程综合、个体化}和\imp{患儿家长共同参与}的原则。\imp{认知行为疗法（CBT）}是当前儿童少年心理治疗的有效方式，其内容主要为识别和纠正患儿在思维、自我形象、态度和信念等方面存在的认知歪曲和错误。对于低年龄的干预对象，可采用低强度的CBT干预为基础治疗，并向其父母提供CBT工作手册，指导父母帮助孩子克服困难。
\end{keybox}

\subsubsection{社交焦虑障碍}

\begin{keybox}
\textbf{社交焦虑障碍（social anxiety disorder）：}又称社交恐惧症，是指在一种或多种社交或公共场合中表现出与环境实际威胁不相称的强烈恐惧和/或焦虑及回避行为。

\vspace{4pt}
\textbf{（1）流行特征：}

\begin{itemize}[leftmargin=*]
    \item 全球社交焦虑障碍患病率在儿童中为4.7\%，在少年中为8.3\%，女童比男童患病率更高，且在少年的性别差异更为明显。
    \item 2020年调查显示，中国在校学生社交焦虑障碍的患病率为0.8\%。
\end{itemize}

\vspace{4pt}
\textbf{（2）主要表现：}

\begin{itemize}[leftmargin=*]
    \item 患儿往往很少认识到自己对社交情境的远离和恐惧是不合理的，可以表现为哭闹、发脾气、依恋他人、畏缩或者不敢在社交情境中讲话。
    \item 回避也是常见的表现形式，且较成年人范围广，轻者表现为过度准备可能需要应付的课堂提问和与同学交流话题，或转移注意力、减少目光接触等；严重者甚至会拒绝上学。
    \item 无论年龄，诊断社交焦虑障碍都要求害怕、焦虑或者回避通常持续至少6个月。
\end{itemize}

\vspace{4pt}
\textbf{（3）预防及干预：}

对患儿的干预主要为心理干预与药物治疗。心理干预建议首选\imp{个别或团体CBT}，其内容主要为心理教育、对恐惧或回避情景的暴露、社交技能训练，以及在真实社交情景中反复训练。干预需考虑与儿童少年的心理发展匹配，且常需父母的配合参与。父母、学校教育方式的调整或阳性强化其社交行为等心理干预方法比药物治疗的效果更好。
\end{keybox}

\subsubsection{分离焦虑障碍（SAD）}

\begin{keybox}
\textbf{分离焦虑障碍（separation anxiety disorder, SAD）：}是儿童少年中常见的焦虑障碍，具体指与特定的依恋对象分离感到显著、过度的恐惧或焦虑的一种精神障碍。

\vspace{4pt}
\textbf{（1）流行特征：}

\begin{itemize}[leftmargin=*]
    \item 不同流行病学调查显示，儿童少年的SAD患病率为1\%$\sim$4\%。
    \item 国内报告6$\sim$16岁在校学生的SAD患病率为0.5\%$\sim$1.0\%。
    \item 儿童少年SAD的危险因素包括焦虑障碍家族史、创伤经历（如依恋对象的疾病或死亡）、生活压力（如父母冲突或离婚）、环境的变化（如搬家或转学）。
\end{itemize}

\vspace{4pt}
\textbf{（2）主要表现：}

\begin{itemize}[leftmargin=*]
    \item 患儿通常对离开家或亲人表现为过度担忧，出现反复的痛苦情绪，达到与其发育水平不符的程度。
    \item 患儿往往对因疾病或灾难失去亲人感到极度焦虑，甚至担心自己会因意外事件（如迷路或被绑架）而与父母分离；也可能会因害怕分离而拒绝离开家，或不愿独自在家或外面过夜。
    \item 其他常见表现还包括频繁做关于分离的噩梦，以及在预期分离时抱怨身体不适，如头痛或胃痛。此外，分离性焦虑障碍也可能与惊恐障碍相关，表现为反复出现突如其来的强烈焦虑和恐惧感。
\end{itemize}

\vspace{4pt}
\textbf{（3）预防及干预：}

如果儿童少年在面对分离场景时的焦虑程度比正常发育阶段的其他人群要严重得多，需引起父母和教师的注意，早期诊断和治疗有助于减轻症状，防止疾病加重。

针对患儿，由精神科医生或其他经过专门培训的合格精神卫生专业人员提供的心理治疗可减少SAD的症状，其中\imp{认知行为治疗（CBT）}被认为是SAD的一线治疗手段。父母也可以学习如何有效地提供情感支持并鼓励孩子形成与年龄相符的独立性。当症状严重或持续存在时，可考虑药物和CBT联合使用。
\end{keybox}

\subsection{抑郁障碍}

\subsubsection{重性抑郁障碍（MDD）}

\begin{keybox}
\textbf{重性抑郁障碍（major depressive disorder, MDD）：}是一种严重且持续的情绪障碍，其核心特征是深度的情绪低落，并伴有兴趣丧失、疲乏、注意力不集中、食欲或睡眠改变、自责感甚至自杀念头等症状，通常持续至少两周或更长时间。

\vspace{4pt}
\textbf{（1）流行特征：}

\begin{itemize}[leftmargin=*]
    \item 中国6$\sim$16岁在校学生MDD的现患率为1.9\%$\sim$2.0\%，其中13$\sim$16岁少年的患病率（2.6\%$\sim$3.0\%）明显高于6$\sim$12岁儿童（1.3\%$\sim$1.5\%）。
    \item MDD在童年期的男女比例相当，青春期后女性多于男性（约为2:1）。
\end{itemize}

\vspace{4pt}
\textbf{（2）主要表现：}

\begin{itemize}[leftmargin=*]
    \item 主要特征包括持续的情绪低落或易怒，对日常活动和兴趣的明显丧失，以及精力不足或过度疲劳。
    \item 常伴有睡眠问题、食欲改变、注意力不集中、学业成绩下降，以及强烈的自责或无价值感。严重时，可能伴随自伤行为或自杀念头。
    \item 可表现出很多焦虑症状（包括恐惧和分离焦虑）和躯体症状（如腹痛、头痛），抑郁经常被躯体症状所掩盖。
\end{itemize}

\vspace{4pt}
\textbf{（3）预防及干预：}

以心理支持为主，常采用行为激活、认知行为疗法、人际心理治疗、解决问题疗法等方法。心理支持可由专业医师或受督导的非专业治疗师教授患儿新的思维方式、应对方式或与他人相处的方式，也可通过自助手册、网站和应用程序获得服务。

在治疗儿童时还须采取针对家庭和学校的并行措施；对于少年，通常采用\imp{心理治疗和抗抑郁药的联合治疗}。
\end{keybox}

\subsubsection{双相障碍（BD）}

\begin{keybox}
\textbf{双相障碍（bipolar disorder, BD）：}也称双相情感障碍，指既有躁狂或轻躁狂发作，又有抑郁发作的一类心境障碍。

\vspace{4pt}
\textbf{（1）流行特征：}

\begin{itemize}[leftmargin=*]
    \item 2019年全球疾病负担研究显示，10$\sim$14岁和15$\sim$19岁儿童少年BD的患病率分别为0.11\%$\sim$0.23\%和0.40\%$\sim$0.80\%。
    \item 对中国6$\sim$16岁在校儿童少年的报告指出，BD的现患率为0.79\%$\sim$0.92\%；在少年期（1.46\%$\sim$1.73\%）明显高于儿童期（0.25\%$\sim$0.36\%）。
\end{itemize}

\vspace{4pt}
\textbf{（2）主要表现（包括躁狂发作和抑郁发作）：}

\begin{enumerate}[leftmargin=*]
    \item \imp{躁狂发作期间}——表现为情绪高涨或易怒、活跃程度显著增加、言语增多且话语连贯性减弱、注意力分散、学习、社交或性行为方面的活动增加，出现冲动或危险行为。躁狂状态可以持续数天至数周。

    \item \imp{抑郁发作期间}——表现为情绪低落、对以前感兴趣的活动失去兴趣或乐趣、活力下降、疲乏或无助感、睡眠障碍（如失眠或嗜睡）、食欲变化以及自杀念头或行为。这种抑郁状态通常持续至少两周，对个体的学习和社交能力产生显著负面影响。

    \item \imp{躯体症状}——患儿还倾向通过行为来表达其情绪，常表述有胃痛、厌食、遗尿、头痛、腹痛、心慌等躯体不适。
\end{enumerate}

\vspace{4pt}
\textbf{（3）预防及干预：}

\begin{enumerate}[leftmargin=*]
    \item 对有家族病史或遗传易感性的儿童少年，应进行心理健康教育和定期筛查，早期发现潜在问题并进行干预，教导儿童识别和管理情绪波动、保持规律的作息时间和健康的生活习惯，帮助其更好地应对压力。
    \item 对患儿的干预方法通常包括\imp{心理治疗}、\imp{家庭治疗}和\imp{药物治疗}。认知行为疗法（CBT）可以帮助儿童识别和改变负面思维模式，培养有效的情绪调节技巧；家庭治疗则强调家庭成员之间的沟通与支持，帮助家庭建立健康的互动模式，以更好地应对疾病带来的挑战。
    \item \imp{学校环境的支持}也十分重要，教育工作者应了解儿童的情况，并提供必要的学习和心理支持，以促进其在校表现和整体发展。
\end{enumerate}
\end{keybox}

\subsection{其他精神障碍}

\subsubsection{强迫性障碍（OCD）}

\begin{keybox}
\textbf{强迫性障碍（obsessive-compulsive disorder, OCD）：}指以强迫思维和/或强迫行为为主要症状，伴有焦虑情绪和适应困难的精神障碍。

\vspace{4pt}
\textbf{（1）流行特征：}

\begin{itemize}[leftmargin=*]
    \item 世界范围内，儿童少年OCD的患病率为1.1\%$\sim$1.8\%。
    \item OCD在中国6$\sim$16岁儿童少年中的患病率为1.0\%$\sim$1.7\%。
    \item 儿童期起病的OCD多见于男性，约30\%的患儿合并终生TD。
\end{itemize}

\vspace{4pt}
\textbf{（2）主要表现：}

\begin{enumerate}[leftmargin=*]
    \item \imp{强迫思维}——是非理性的、不由自主重复出现的思想、观念、表象、意念、冲动等，这些思维是不愉快的、非自愿的，而且在绝大多数个体身上导致显著的痛苦或焦虑。个体企图忽略或压抑这些强迫思维，或用其他想法、行动（例如，执行一个强迫行为）来中和它们。

    \item \imp{强迫行为}——是反复的行为或精神活动，如清洗、检查、数数和反复默诵，患儿可能感到被强迫思维驱使，或认为必须机械地遵守规则，因而不得不去执行。

    \item 绝大多数有OCD的个体既有强迫思维又有强迫行为，强迫行为通常作为对强迫思维的反应而实施，目标是减少被强迫思维所激发的痛苦，或防止所担心的事件发生。
\end{enumerate}

\vspace{4pt}
\textbf{（3）预防及干预：}

预防OCD的策略更加注重识别和管理过度的焦虑、强迫性思维、以及不合常规的反复行为，如过度清洁、检查或重复特定动作，这些行为一旦影响到日常生活，就需要引起警惕。

OCD的干预须\imp{心理干预和药物治疗相结合}。心理干预可采用心理教育、支持性心理治疗、家庭治疗、认知行为治疗，也可选择森田疗法、生物反馈疗法、音乐疗法等，可根据患儿具体症状灵活选用。
\end{keybox}

\subsubsection{品行障碍（CD）}

\begin{keybox}
\textbf{品行障碍（conduct disorder, CD）：}指儿童少年（18周岁以下）反复出现持久的、违反其年龄段相应行为规范，侵犯他人或公共利益的行为障碍，包括反社会性、攻击性或对立违抗行为。

CD往往继发于童年期的\imp{对立违抗障碍（oppositional defiant disorder, ODD）}，表现为持久性的违抗、敌意、不服从、易激惹、挑衅和破坏行为等，且具有冲动性，少数严重者可进一步发展为少年违法犯罪，即参与违反法律或与社会规则相悖的违法行为。

\vspace{4pt}
\textbf{（1）流行特征：}

\begin{itemize}[leftmargin=*]
    \item CD在青春早期高发，男性显著多于女性[（4$\sim$12）:1]。在全球范围内，CD患病率约为8\%。
    \item CD在中国6$\sim$16岁在校学生中的患病率约为1.9\%，10$\sim$14岁组的患病率最高。
\end{itemize}

\vspace{4pt}
\textbf{（2）主要表现：}

\begin{itemize}[leftmargin=*]
    \item 患有CD的儿童少年对他人的感受缺乏敏感性，有时会将他人的行为误认为是威胁，进而表现出攻击性行为。
    \item 患儿常易激惹、固执、爱寻衅、恶作剧、侵犯和攻击他人、违拗对抗、故意破坏财物、虐待动物、撒谎欺诈、持续而顽固的偷窃、厌学逃学、离家出走等。
    \item 严重时可进一步发展为暴力伤人、性侵犯、诈骗、偷盗、吸毒等违法犯罪行为，此时多伴有反社会人格障碍。
    \item 这些行为均违反了与年龄相适应的社会行为规范和道德准则，危害社会安定及他人、公众利益，而且对障碍者自身的学习、生活和社会化功能也造成严重不良影响。
\end{itemize}

\vspace{4pt}
\textbf{（3）预防及干预：}

品行障碍比较顽固，矫治十分困难，因此\imp{及早发现}、\imp{及早采取预防性干预}极为重要。
\end{keybox}

\subsubsection{神经性厌食症（AN）}

\begin{keybox}
\textbf{神经性厌食症（anorexia nervosa, AN）：}是以持续性的能量摄取限制、强烈害怕体质量增加或变胖或持续性妨碍体质量增加的行为、对自我的体质量或体形产生感知紊乱为临床特征的一类进食障碍。

\vspace{4pt}
\textbf{（1）流行特征：}

\begin{itemize}[leftmargin=*]
    \item AN发病年龄通常较早，为13$\sim$20岁，发病的两个高峰年龄分别是13$\sim$14岁和17$\sim$18岁。
    \item 男女发病率存在较大差异[1:（11$\sim$14）]。
    \item 中国6$\sim$16岁在校学生的AN现患率为0.1\%$\sim$0.2\%。
\end{itemize}

\vspace{4pt}
\textbf{（2）主要表现：}

\begin{enumerate}[leftmargin=*]
    \item \imp{限制饮食摄入}——在行为上，患者常刻意减少饮食的摄入量，并伴有过度运动以增加消耗。
    \item \imp{使用其他减重方法}——例如使用泻药、减肥辅助剂或进食后催吐。
    \item \imp{体像障碍}——在心理上"迷恋"低体重，存在体像障碍，对自身体形的感知异常，如明显已经很消瘦了，仍觉得自己很胖。
    \item \imp{闭经}——对于已经开始正常月经的女孩来说，患有AN的女性都会停止月经，关于腹胀、胃痉挛和便秘的报告也很常见。
\end{enumerate}

\vspace{4pt}
\textbf{（3）预防及干预：}

预防儿童少年AN的重点是建立健康的饮食观念和积极的身体形象。家庭和学校应共同关注儿童的心理发展，增强自尊与自信，避免将体重和外貌作为评价标准。社会和学校需引导健康的审美观念，提供心理支持，创造多样化文化环境，预防进食障碍的发生。

治疗的核心目标是\imp{恢复体质量}，关键在于早期诊断和及时营养重建。饮食监管和禁止暴食与呕吐行为是主要的治疗手段。由于AN常伴随心境障碍、焦虑障碍和强迫症等精神障碍，可根据症状进行药物干预。为患者及家庭提供心理教育，建立治疗联盟，并通过系统的心理行为干预实施全程管理，以提高治疗效果和促进整体恢复。
\end{keybox}

\newpage
